\paragraph{Exercise 1.} If $V\subset U$ is affine open, then $V$ is open in $X$ and its restricted sheaf from $U$ equals its restriction from $X$.
\paragraph{Exercise 2.} It is $\Spec(\mathbb Z/n\mathbb Z)$ with support the primes $(p)$ dividing $n$ (empty when $n$ is a unit).
\paragraph{Exercise 3.} All have support $\{(x)\}$, but the quotients have nilpotents of different orders.
\paragraph{Exercise 4.} The subscheme is $\Spec k$, so its local ring at the unique point is $k$.
\paragraph{Exercise 5.} For every prime $\mathfrak p$, $I\subseteq\mathfrak p$ iff $\sqrt I\subseteq\mathfrak p$.
\paragraph{Exercise 6.} $(x)$ and $(x^2)$ have radical $(x)$; the quotients are $k$ and $k[x]/(x^2)$.
\paragraph{Exercise 7.} It is $\Spec k[y,y^{-1}]$, since $k[x,y]_y/(x)\cong k[y,y^{-1}]$.
\paragraph{Exercise 8.} Emptiness means $A_f/IA_f=0$, hence $1\in IA_f$; clearing denominators gives $f^n\in I$.
\paragraph{Exercise 9.} If $(R,\mathfrak m)$ is local and $I
e R$, then $R/I$ has unique maximal ideal $\mathfrak m/I$.
\paragraph{Exercise 10.} For $A\twoheadrightarrow A/I$ and $\mathfrak p\supset I$, both residue fields equal $\operatorname{Frac}(A/\mathfrak p)$ in the usual prime-quotient sense.
\paragraph{Exercise 11.} If $X$ is reduced, each affine coordinate ring has nilradical zero, so quotienting by it changes nothing.
\paragraph{Exercise 12.} $\sqrt{(x^2y^3)}=(xy)$, so the reduction is $\Spec k[x,y]/(xy)$.
\paragraph{Exercise 13.} After one reduction all affine coordinate rings are reduced, so a second reduction is the identity.
\paragraph{Exercise 14.} A map to a reduced ring sends every nilpotent to zero, so it kills $\sqrt{(0)}$.
\paragraph{Exercise 15.} They are equal as subsets: $V(I+J)=V(I)\cap V(J)$.
\paragraph{Exercise 16.} They are equal as subsets: $V(I\cap J)=V(I)\cup V(J)$.
\paragraph{Exercise 17.} The sum $(x)+(y)=(x,y)$ gives the reduced origin $\Spec k$.
\paragraph{Exercise 18.} The intersection $(x)\cap(y)=(xy)$ gives $\Spec k[x,y]/(xy)$.
\paragraph{Exercise 19.} The open immersion $D(t)\hookrightarrow\mathbb A^1_k$ is an immersion, and its image $D(t)$ is not closed in $\mathbb A^1_k$.
\paragraph{Exercise 20.} A homeomorphism onto an open subset is injective.
\paragraph{Exercise 21.} A homeomorphism onto a closed subset is injective.
\paragraph{Exercise 22.} By the two definitions, the image subset is both open and closed and the sheaf identifications agree.
\paragraph{Exercise 23.} The quotient $A/I\twoheadrightarrow A/J$ has kernel $J/I$.
\paragraph{Exercise 24.} Both rings receive $A$ with $I$ killed and $f$ inverted, so the universal property gives a canonical isomorphism.
\paragraph{Exercise 25.} The open complement is an open subset, whose structure is uniquely $\OO_X$ restricted to it; no ideal choice is needed.
\paragraph{Exercise 26.} $\Spec k[\varepsilon]/(\varepsilon^2)$ and $\Spec k$ have the same one-point topology, with the latter the reduction.
